\documentclass[a4paper,twoside,12pt]{book}

%\usepackage{CJK}
%\usepackage{xeCJK}
\usepackage{amssymb}
\usepackage{fontspec}
\setmainfont{cwTeX 明體}
\newfontfamily\mylt{jsMath-cmti10}
\newfontfamily\mysm{Courier New}
%\setmainfont[ItalicFont]{jsMath-cmti10}

%\setmainfont{cwTeX 粗黑體}

%\newfontfamily\cwbm{cwTeX 粗黑體}
%\newfontfamily\cwbm{cwTeX 楷書}
%\setCJKmainfont{cwTeX 明體}

\usepackage{multicol}
\usepackage{graphicx}
\usepackage{bm, amsmath, url}
\usepackage{amsfonts, keystroke}
\usepackage{appendix, enumitem, listings}
\usepackage{fancyvrb}
\usepackage[Lenny]{fncychap}
\usepackage{fullpage} 
%minitoc}
\usepackage{xcolor}
\definecolor{bgcolor}{rgb}{0.9,0.9,0.9}
%定义背景颜色
\lstloadlanguages{[LaTeX]TeX}

%\setcounter{parttocdept}{3}
%\setlength{\ptcindent}{0pt}
%\renewcommand{\mtcfont}{\footnotesize\cwbm} 
%\renewcommand{\mtcSfont}{\small}

\newenvironment{latexample}[1][language={[LaTeX]TeX}]
{\lstset{backgroundcolor=\color{bgcolor},
    keywordstyle=\color[rgb]{0,0,1},
    commentstyle=\color[rgb]{0.133,0.545,0.133},
    stringstyle=\color[rgb]{0.627,0.126,0.941},
    breaklines=true,
    prebreak = \raisebox{0ex}[0ex][0ex]{\ensuremath{\hookleftarrow}},
    frame=single,
    language={[LaTeX]TeX},
    basicstyle=\footnotesize\ttfamily, #1}
  \VerbatimEnvironment\begin{VerbatimOut}{latexample.verb.out}}
  {\end{VerbatimOut}\noindent
  \begin{minipage}{0.5\linewidth}
    \lstinputlisting[]{latexample.verb.out}%
  \end{minipage}\qquad
  \begin{minipage}{0.4\linewidth}
    \input{latexample.verb.out}
  \end{minipage}\\
}

\newenvironment{latexamplefull}[1][language={[LaTeX]TeX}]
{\lstset{backgroundcolor=\color{bgcolor},
    keywordstyle=\color[rgb]{0,0,1},
    commentstyle=\color[rgb]{0.133,0.545,0.133},
    stringstyle=\color[rgb]{0.627,0.126,0.941},
    breaklines=true,
    prebreak = \raisebox{0ex}[0ex][0ex]{\ensuremath{\hookleftarrow}},
    frame=single,
    language={[LaTeX]TeX},
    basicstyle=\footnotesize\ttfamily, #1}
  \VerbatimEnvironment\begin{VerbatimOut}{latexample.verb.out}}
  {\end{VerbatimOut}\noindent
  \centering\begin{minipage}{0.8\linewidth}
    \lstinputlisting[]{latexample.verb.out}%
  \end{minipage}\qquad
  \\
}
%\renewcommand{\chaptername}{Chapter~\arabic{chapter}}
%\renewcommand{\thechapter}{\cwbm}
%\renewcommand{\thesection}{\cwbm{\arabic{chapter}.\arabic{section}}}
%\renewcommand{\thesubsection}{\cwbm{\arabic{chapter}.\arabic{section}.\arabic{subsection}}}
\VerbatimFootnotes


\addtolength{\headsep}{+1cm}
\addtolength{\oddsidemargin}{+0.5cm}
\addtolength{\evensidemargin}{-0.5cm}

%%\voffset=-0.5cm
%\hoffset=-0.5cm
%%\setlength{\topmargin}{-0.15in}
%\setlength{\topmargin}{0.5cm}
%\setlength{\headheight}{0in}
%\setlength{\headsep}{0.5in}
%\setlength{\footskip}{1in}
%\setlength{\textwidth}{17cm}
%\setlength{\textheight}{23cm}
%\setlength{\marginparsep}{0in}
%\setlength{\marginparwidth}{0in}
%%\setlength{\oddsidemargin}{0in}
%%\setlength{\evensidemargin}{-0in}


\usepackage{fancyhdr}
\usepackage{booktabs}
%\usepackage[cam,a4,center]{crop}
%\usepackage{geometry}
\pagestyle{fancy}

\renewcommand{\headrulewidth}{0.4pt}
%\renewcommand{\headrulelength}{\textwidth}
\renewcommand{\sectionmark}[1]{\markright{\sl{\thesection.\ #1}}}
\renewcommand{\leftmark}{{\sl Chapter}.\thechapter }
%\lhead{\footnotesize\rightmark}
%\rhead{\thepage}
\fancyfoot{}
\fancyhead[RE]{\footnotesize\leftmark}
\fancyhead[LO]{\footnotesize\rightmark}
\fancyhead[LE,RO]{\thepage}



\XeTeXlinebreaklocale "zh"
\XeTeXlinebreakskip = 0pt plus 1pt
\setlength{\parskip}{0.3cm}

%%%%%%%%%%%%%%
\newcommand{\engline}[1]{\renewcommand\baselinestretch{1}\selectfont #1\par~\\[-36pt]}
\newcommand{\myepigraph}[2]{\noindent~~~~~~~~~~~~~~~~~~~\begin{minipage}{0.8\textwidth}\setlength{\parsep}{0pt} \mysm\begin{small}\engline{#1}
\end{small}\\[-10pt]\hline\end{minipage}\end{center}\begin{flushright}\mysm \begin{small}~\\[-24pt]#2\end{small}~~~~~\end{flushright}\\~\\[-24pt]}



%%%%%%%%%%%%%%
\newcommand{\btn}[1]{
\fboxrule=0.7pt
\fboxsep=1.2pt
\fbox{#1}}
\newcommand{\btnn}{\Enter}
%\newcommand{\btnn}{\fboxrule=0.7pt
%\fboxsep=1.2pt
%~~~\fbox{\textsc{Enter}}}

\newcommand{\ex}{\addtocounter{subsection}{+1}
\subsubsection{Ex~\underline{\thesubsection}}}
\newcommand{\ans}{
\subsubsection*{\underline{Ans:}}}
\newcommand{\resl}[1]{\textcolor{blue}{#1}}
\newcommand{\reslc}[1]{\textcolor{teal!65!black}{\tt{//}~\textrm{#1}}}
\newcommand{\reslp}[1]{\textcolor{blue}{\tt{\%}~\textrm{#1}}}
\newcommand{\resle}[1]{\textcolor{blue}{\textrm{#1}}}
\newcommand{\scicode}[1]{\bigskip\noindent\fbox{\begin{minipage}{0.95\textwidth}~\begin{minipage}{0.95\textwidth}~\\[-5pt]#1\end{minipage}\\\smallskip\end{minipage}}\bigskip\vline\vline\vline}

\newcommand{\norm}[1]{\|#1\|}
%\newcommand{\sdelta}{{\scriptstyle\Delta}}
\newcommand{\sdelta}{\Delta}
\newcommand{\bra}[1]{\langle#1|}
\newcommand{\ket}[1]{|#1\rangle }
\newcommand{\bket}[1]{\Big |#1\Big\rangle }
\newcommand{\braket}[2]{\langle#1|#2\rangle }
\newcommand{\proj}[2]{|#1\rangle\langle#2|}
\newcommand{\expect}[2]{\langle #1\rangle_{#2}}
\newcommand{\expecto}[1]{\langle #1\rangle}
\newcommand{\Kappa}[0]{\bm{\mathcal{K}}}
\newcommand{\parity}[0]{\bm{\mathcal{P}}}
\newcommand{\parityl}[0]{\bm{p}}
\newcommand{\charge}[0]{\bm{\mathcal{C}}}
%\DefineShortVerb{\|\#}
%%%%%%%%%%%%%%%

%\setitemize{noitemsep}
%\setenumerate{noitemsep}
\setitemize{itemsep=-5pt}
\setenumerate{itemsep=-5pt}
\linespread{1.35}


\begin{document}
\begin{titlepage}
\title{量子力學的數學結構}
\author{Kow Lung Chang 張國龍\\
Physics Department, National Taiwan University\\~\\~\\}
\maketitle
\end{titlepage}


%
%\begin{titlepage}
%
%\title{量子力學的數學結構}
%\author{Kow Lung Chang~張國龍\\\\
%Physics Department, National Taiwan University} 
%
%\date
%\maketitle
%\end{titlepage}

\pagestyle{empty}
\frontmatter


~\\~\\~\\~\\~\\~\\~\\
\begin{center}
{\LARGE \sl TO FEZA AND SUHA}
\end{center}


~~~
\newpage

\setcounter{page}{7}
\pagestyle{plain}
\begin{raggedleft} {\LARGE 引言}\end{raggedleft}\\\\

研究不同的學門所使用的數學工具會有所差異; 在古典力學的領域裡，我們需要具備變分學（variational calculus）的知識，才能處理一群（或一個）質點的動力方程式。在量子力學裡，也一樣需要一些不同的數學工具。我們需要線性代數（linear algebra）, 線性向量空間（linear vector space）的知識，更重要的是對 Hilbert 空間 （Hilbert space）的了解和熟稔，才能完全領會量子力學最基本的四大公理（the four postuates of quantum mechanics）。

不同的探討領域，也還需要不同的數學工具。量子力學系統中，很重要的一個特徵是系統的對稱性（symmetry property）。當我們研究系統的對稱轉換（symmetry transformation）時，就需要用到連續群（continuous group）或李群（Lie group）的理論及其 group representations。

從我過去講授量子力學課程的經驗裡，時常遇到一個由學生提出的共同問題：到底我們要具備多少數學的背景知識，才能有信心、有效率地來研修這門課程。當然這個問題難給予統一的答案，因為學生當中所具備的數學背景頗為懸殊，一般大學部標準的應用數學課程，一定包括線性代數等範圍，有的學生也會到數學系選修群論（group theory）等課程。當我無法在講授量子力學的同時解決擁有不同數學背景學生們遇到的困難時，我只能數度修訂我的量力講義，當需使用到特殊的數學工具時，我將這些資料先行介紹給學生，節省他們重新摸索的時間。

這本書首先回顧線性向量空間，並以其數學術語來闡述量子力學的基本定律。第一章除了介紹量子力學前三大定律外，也討論在量子力學中不可或缺，線性算子（linear operators）和 dynamical observables 之間的關係。

量子力學當然不只是解某一特定的系統的 Schr\"{o}dinger 方程式，但是傳統上它卻一再用來作為系統量化的實例： 如氫原子光譜，或者它的角動量。本書第二章藉由將前一章抽象的量子理論數學架構轉化為讀者熟悉的 q-representation 與 p-representation ，同時輔以經典的簡諧振子（harmonic oscillator）為例，說明 N-representation 與 q-representation 之間的關係。

很多作者試圖以波包（wave packet）的方法來推導 Schr\"{o}dinger 方程式。嚴格來說，Schr\"{o}dinger 方程式絕對不是波包運動時省略某些項而獲得的近似結果。Schr\"{o}dinger 方程式是一個 exact equation，並且是量子力學中唯一的 dynammical postulate。正如同古典力學中，牛頓第二運動定律：$\vec{f}=m\vec{a}$ 是 dynamical law of motion 而不是可從推導去獲得一般，Schr\"{o}dinger 方程式這個動力學上的假設決定了一個量子態（quantum state） 隨時間的演化（time evolution）。在第二章結束以前，我們將示範在古典理論的極限下，即 $\hbar\rightarrow 0$ 時，古典力學的運動方程式確實可以從量子系統時域演化所定義的傳播子（propagator）之極大值再次獲得。

從第三章到最後一章，探討的重心完全集中於量子系統的對稱性。在第三章中，我們分別探討連續群（continuous group）與離散群（discrete group）; 介紹李群（Lie group）與李代數（Lie algebra），並且深入分析 半純李群（semisimple Lie group）及 半純李群代數（semisimple Lie algebra） 之標準形（standard form）。至於離散群的對稱性則僅限於 PCT 的分析。

第四章則由 O(3) 及 SU(2) 對稱群來探討角動量，並介紹 O(3) 的不可約表示（	irreducible representation ）與球諧函數（spherical harmonics）之間的關聯。最後也介紹 dynamical symmetry 在 Coulomb 力作用之系統中的應用來做總結。

最後一章從時空的對稱轉換來建構 Lorentz 群，並且從 Lorentz 群的有限維表示（finite dimensional representation）來進行具有不同內稟自旋（intrinsic spin）粒子的分類。我們特別以 Dirac 粒子為例詳細地探討 Dirac 方程式及其相關的 PCT 轉換的特徵。

本書的專門術語將儘量依照國家教育研究院公告之學術名詞對照表給予中文譯名，並在該名詞首次出現時同時列出英文原名。少部份無適當譯名之術語則照列原文。另外也可以從附錄的中英名詞對照表來檢索。



\begin{flushleft}
Taipei, Taiwan\\ 
September, 2011\hfill Kow Lung Chang
\end{flushleft}

\newpage

\setcounter{page}{9}
\pagestyle{fancy}
\fancyhead[RE]{\footnotesize\textit{Contents}}
\fancyhead[LO]{\footnotesize\textit{Contents}}
\pretolerance=10000
\begin{small}
\tableofcontents
\end{small}
\pretolerance=100

\mainmatter
\fancyhead[RE]{\footnotesize\leftmark}
\fancyhead[LO]{\footnotesize\rightmark}
%\thispagestyle{headings}
\pagestyle{fancy}
%\chapter{量子力學的基本定律及相關定理}

量子力學的理論架構建立在四個基本的定律上，正如同古典力學建構於牛頓三個運動定律形成之基石。不同物理領域的研究者會使用形色各異的數學工具，譬如牛頓為了精確描述物體的運動現象而發明 fluxion，即今日所稱的微分。幾世紀以來，數學家已經發展出古典力學以外形形色色的數學分支，並逐漸演變成各自蓬勃發展的子領域，而物理學家從中選用適合特定物理系統的數學作為理論推演與計算的工具。我們將應用 Hilbert 空間與線性算子（linear oprators）等數學工具來分析量子力學的基本定律。

為了讀者方便，本章開頭會以幾個章節來回顧線性向量空間（linear vector space）、 Hilbert 空間與線性算子（linear oprators）等背景資料，讓不熟悉這些數學的讀者可以省去自我摸索的時間。自信已具有這方面知識的讀者可以考慮跳過這些章節，直接從第 5 節開始閱讀。

\section{向量空間}

我們用符號 $\mathcal{V}$ 代表一個向量空間（vector space），它是由無數多個元素（element）所構成之集合（set），這些元素稱為向量（vectors），平常以斜體的希臘字母 $\psi, \varphi, \chi, \ldots$ 表示。向量空間的元素滿足下述兩種操作（operation）：

\begin{itemize}\itemsep=-2pt
\item 可加性 （An operation of addition）：對於任意集合中任意一對向量 $\psi$ 與 $\phi$，可對應到一個新的向量 $ \psi +\phi \in \mathcal{V} $, 我們稱它為 $\psi$ 與 $\phi$ 的和。

\item 純量積（An operation of scalar multiplication）：對應於 $ \psi $ 與一個純數 $ a $，則 $ a\psi $ 也是該集合中的一個向量。假如 $a,b$ 皆為純數，$ \psi , \phi $ 與 $ \chi $ 為線性空間的向量，則這些向量還必須滿足下面 8 個條件：
\addtocounter{equation}{1}
\begin{align}
~&\psi +\phi=\phi +\psi,\tag{\arabic{chapter}.\arabic{equation}a}\\
~&\psi +(\phi +\chi)=(\psi +\phi )+\chi,\tag{\arabic{chapter}.\arabic{equation}b}\\
~&\psi +\mathit{0}=\psi, ~~\mathit{0}~\mbox{is null vector},\tag{\arabic{chapter}.\arabic{equation}c}~~~~~~~~~~~\\
~&a(\psi +\phi)=a\psi +a\phi,\tag{\arabic{chapter}.\arabic{equation}d} \\
~&(a+b)\psi=a\psi+b\psi,\tag{\arabic{chapter}.\arabic{equation}e} \\
~&a(b\psi)=(ab)\psi,\tag{\arabic{chapter}.\arabic{equation}f}\\ 
~&1\cdot\psi=\psi,\tag{\arabic{chapter}.\arabic{equation}g}\\
~&0\cdot \psi=\mathit{0},\tag{\arabic{chapter}.\arabic{equation}h} 
\end{align}
\end{itemize}

當 $a,b$ 皆為實數時，我們將該向量空間稱為實向量空間（real vector space），並以 $ \mathcal{V}_r $ 表示。如果 $a,b$ 皆為複數（complex number），則該該向量空間稱複向量空間（complex vector space）$ \mathcal{V}_c $。

\subsubsection*{範例一}
我們以讀者熟悉的 $ \mathcal{R}^n $ 空間，即 $ n $ 維實數歐式空間（$ n $-dimensional Euclidean space）為例，$ \mathcal{R}^n $ 空間的向量 $ \psi $ 與 $ \phi $ 可表示成 $ \psi =(x_1, x_2,\ldots,$ $x_i, \ldots, x_n) $ 與 $ \phi =(y_1,y_2,\ldots,y_i,\ldots,y_n) $， 其中 $ x_i $ 和 $ y_i $ ($ i=1,2,\ldots,$ $ n $) 皆為實數，且 $ \psi $ 與 $ \phi $ 之和可表以 $ (x_1+y_1,x_2+y_2,\ldots,x_i+y_i,\ldots,x_n+y_n) $。$\psi$ 乘以實數 $a$ 則表以 $ a\psi =a(x_1,x_2,\ldots,x_i,\ldots,x_n)= (ax_1,ax_2,\ldots,ax_i,\ldots$ $ ,ax_n) $。 若 $ a $ 與 $ x_i $ 或 $ y_i $ 皆以複數代之，則該向量空間稱為 $n$ 維複數歐式空間（$ n $-dimensional complex Euclidean space）並以 $ \mathcal{C}^n$ 表示。讀者可以自行檢驗不論實或複複數歐式空間中的向量 $\psi$ 或 $ \phi $ 皆滿足 Eqs.~(1.1a) 至 (1.1h) 的條件。

\subsubsection*{範例二}

在 $ a\leqslant x\leqslant b$ 的範圍內，如果 $ f(x) $ 代表所有連續函數的集合，那麼這些函數也可以形成一個向量空間，因為在可加性與純量積兩種操作下，連續函數也會滿足 Eqs.~(1.1a) 至 (1.1h)之條件。

\subsection{線性獨立與線性相依}
如果 $ \{\psi_1,\psi_2,\ldots, \psi_m\} $ 代表向量空間中的 $ m $ 向量，並且這 $ m $ 個向量建構如下所示之線性組合

\begin{align}
\sum_{i=1}^m a_i\psi_i.
\end{align}

\noindent 當然這 $ m $ 個向量構成的線性組合也是該空間中的向量。如果在所有係數 $ a_i $ 皆為 0 的時候，即 $ a_i =0$， $ i=1,2,\ldots, m $ 才能使該線性組合為零向量（null vector），換言之  

\begin{align}
\sum_{i=1}^m a_i\psi_i=0
\end{align}

只有所有係數皆為 0 時才成立，那麼這組 $ m $ 個向量就稱為線性獨立（linear independent）。相反地，如果其中至少有一個係數不為 0，例如 $ a_l \neq 0$，也可以使 $ \sum_{i=1}^ma_i\psi_i=\mathit{0} $，則這組 $ m $ 個向量稱為線性相依（linearly dependent），也就是因為 $a_j\neq 0$ ，所以 $\psi_j$ 可以寫成

\begin{align*}
\psi_j=\frac{-1}{a_j}\sum_{i\neq j}^m a_i\psi_i.
\end{align*}

這個式子表示 $\psi_j$ 可以表示成其他 $m-1$ 個向量的線性組合，換言之相依於此 $m-1$ 個向量。

\subsection{維數與基底}

在向量空間 $ \mathcal{V} $ 中，所能找到最大（或最多）的線性獨立向量數目即代表該向量空間的維數（dimension）。如果 $ \mathcal{V}_n $ 代表 $n$ 維的向量空間，那麼任何一組 $n$ 個線性獨立的向量就可成為該向量空間的基底（basis）。

\subsubsection*{範例三}

以我們熟悉的三維空間為例，一般它的基底都是選取互相正交（orthogonal）的三個單位向量，即 $\hat{i}=(1,0,0), \hat{j}=(0,1,0), \hat{k}=(0,0,1)$ 。然而基底在向量空間中的定義並不包含正交性的要求，也不需要是單位向量，只要任何 3 個向量滿足線性獨立的要求：也就是其線性組合只有在全係數為零時才等於零向量的條件，則該組 3 個向量即可成為基底，如

\begin{align*}
x_1=(a,0,0),~~~x_2=(a,b,0),~~~x_3=(a,b,c)
\end{align*}

這 3 個向量即構成三維空間的基底，當然物理學家選取正交單位向量作為基底在實際計算上方便許多。在往後數個章節我們會介紹將一般基底轉化成互相正交且屬於單位向量基底的方法。


\section{內積}

向量空間中的內積（inner product）有時稱為純量積（scalar product），藉由操作一對前後有序的向量製造出的一個純數，它可以是實數或者是複數。如果以 $ \psi $ 與 $ \phi $ 代表操作的向量對，那麼我們以符號 $ (\psi,\phi) $ 來標示內積結果。內積的定義中要求這對向量保持秩序，所以 $ (\psi,\phi) $ 並不一定等於 $ (\phi,\psi) $。假如內積得到的純數滿足下列四個條件：

\addtocounter{equation}{1}
\begin{align}
&(\psi,\phi+\chi)=(\psi,\phi)+(\psi,\chi), \tag{\arabic{chapter}.\arabic{equation}a} \\
&(\psi,a\phi)=a(\psi,\phi),\tag{\arabic{chapter}.\arabic{equation}b}\\ 
&(\psi,\phi)=(\phi,\psi)^*,\tag{\arabic{chapter}.\arabic{equation}c}\\ 
&(\psi,\psi)\geqslant 0,~~~\mbox{只有當}~\psi=0 ~\mbox{時 } (\psi,\psi)=0 .\tag{\arabic{chapter}.\arabic{equation}d}
\end{align}

那麼 $ (\psi,\phi) $ 就定義為該向量空間，由有序的一對向量 $ \psi $ 及 $ \phi $ 對應的內積。由前述條件我們馬上可以確定 $(\psi,\psi)$ 一定是實數。如果 $ \psi $ 及 $ \phi $ 都不是零向量，但是 $ (\psi, \phi)=(\phi, \psi)=0 $，那麼我們稱 $ \psi $ 與 $ \phi $ （或 $ \phi $ 與 $ \psi $）互相正交（orthogonal）。

\subsubsection*{範例四}

如果 $ \psi=(x_1,x_2,\ldots,$ $ x_n) $ 和 $ \phi=(y_1,y_2,\ldots, y_n) $ 代表 $ \mathcal{C}^n $ 空間的一對向量，那麼它們的內積可以表示成

\begin{align}
\displaystyle(\psi,\phi)=\sum_{i=1}^n x_i^*y_i=x_1^*y_1+x_2^*y_2+\cdots +x_n^*y_n.
\end{align}

讀者不難驗證上式滿足內積操作要求的四個條件。

\subsubsection*{範例五}

在範例二中我們曾確認在 $ a\leqslant x\leqslant b $ 範圍中的所有連續函數 $ f(x) $ 構成一個向量空間。在該向量空間中，如果 $ f(x) $ 及 $ g(x) $ 代表有序的一對函數，則它們的內積可以表為

\begin{align}
(f(x),g(x))=\int_a^bf(x)^*g(x)dx.
\end{align}

上式的 $ f(x) $ 和 $ g(x) $ 可以是複函數，即 $f(x)=f_1(x)+if_2(x)$。如果 $ |f(x)|^2 $ 和 $ |g(x)|^2 $ 都為有限值（finite），則其各自的內積 $ (f(x),f(x)) $ 和 $ (g(x),g(x)) $ 也都是有限值，則其所屬之向量空間稱為 $ \mathcal{L}^2(a,b)$ 空間。

\subsection{Schwarz 不等式}


探討過向量的內積操作之後，我們可以用這些性質進一步推導在量子力學中常使用的 Schwarz 不等式。對應於向量 $\psi$ 及 $\phi$ ， Schwarz 不等式常表以

\begin{align}
|(\psi,\phi)|=\sqrt{(\psi,\phi)(\phi,\psi)}\leqslant\sqrt{(\psi,\psi)}\sqrt{(\phi,\phi)}.
\end{align}

證明如下：

令 $ \alpha =\xi+i\eta $ 為一複數，則 $\chi=\psi+\alpha\varphi$ 的內積

\begin{align*}
(\chi, \chi)&=(\psi+\alpha\varphi, \psi+\alpha\varphi)\\
&=\xi\{(\psi, \varphi)+(\varphi, \psi)\}+i\eta\{(\psi, \varphi)-(\varphi, \psi)\}+(\xi^2+\eta^2)(\varphi, \varphi)\geqslant 0.
\end{align*}

$(\chi, \chi)\geqslant 0$ 對於任何的 $\xi$ 及 $\eta$ 值皆成立，換句話說，如果將 $(\chi, \chi)$ 視為 $\xi$ 及 $\eta$ 的函數 $f(\xi, \eta)$，$(\chi, \chi)$ 的極小值也大於等於零，即 $(\chi, \chi)_{\mbox{min}}\geqslant 0$。若 $(\chi, \chi)=f(\xi, \eta$ 在 $\xi=\xi_0, \eta=\eta_0$ 時有極小值


\begin{align}
\frac{\partial f(\xi, \eta)}{\partial \xi}\Bigg|_{\xi_0,\eta_0}=\frac{\partial f(\xi, \eta)}{\partial \eta}\Bigg|_{\xi_0,\eta_0}=0,
\end{align}

實際計算可以得到

\begin{align*}
\frac{\partial f(\xi, \eta)}{\partial \xi}\Bigg|_{\xi_0,\eta_0}&=2(\varphi, \varphi)\xi+(\psi, \varphi)+(\varphi, \psi)=0,\\
\frac{\partial f(\xi, \eta)}{\partial \eta}\Bigg|_{\xi_0,\eta_0}&=2(\varphi, \varphi)\eta+i(\psi, \varphi)-i(\varphi, \psi)=0.
\end{align*}

整理後得到

\begin{align}
\xi_0=-\frac{1}{2}\frac{(\psi,\phi)+(\phi,\psi)}{(\phi,\phi)},~~~~
\eta_0=-\frac{i}{2}\frac{(\psi,\phi)-(\phi,\psi)}{(\phi,\phi)}.
\end{align}

因此

\begin{align}
(\chi, \chi)_{\mbox{min}}=f(\xi_0,\eta_0)=(\psi,\psi)-\frac{(\psi,\phi)(\phi,\psi)}{(\phi,\phi)}\geqslant 0,
\end{align}

移項後即常見的 Schwarz 不等式 $|(\psi,\varphi)|^2\leqslant (\psi, \psi)(\varphi, \varphi)$。

內積的特性也可以將向量空間中一組既不是單位長度也不正交的基底轉化為另一組正交且具有單位長度的新基底，這個轉化過程稱為 Gram-Schmidt 正交化。

\subsection{Gram-Schmidt 正交化}

$\{ \psi_1, \psi_2,$ $\ldots , \psi_n \}$  代表 $n$ 維向量空間中的一組基底，一般而言 $(\psi_i,\psi_j)\neq 0, i\neq j $。也就是說，$\{ \psi_1, \psi_2,$ $\ldots , \psi_n \}$ 只具備線性獨立的性質，但不一定互相正交。然而我們可以根據這一組已有的基底，分兩個步驟來建構一組互相正交且為單位長度的正則基底（orthonomal basis）。

首先，我們先尋找 $n$ 個線性獨立的向量 $\psi_i (i=1,2,\ldots, n)$ ，令 $ \psi_1'=\psi_1 $，$ \psi_2'=\psi_2+\alpha\psi_1' $。為了滿足 $ \psi_1' $ 與 $ \psi_2' $ 互相正交的條件，我們可以適當地選擇 $ \alpha $ 使 $ (\psi_2',\psi_1')=0 $，即


\begin{align}
(\psi_2',\psi_1')=(\psi_2,\psi_1')+\alpha^*(\psi_1',\psi_1')=0,
\end{align}

因此得到 $ \alpha=-(\psi_2,\psi_1')^*/(\psi_1',\psi_1')=-(\psi_1',\psi_2)/(\psi_1',\psi_1') $，且

\begin{align}
\psi_2'=\psi_2-\psi_1'\frac{(\psi_1',\psi_2)}{(\psi_1',\psi_1')}.
\end{align} 


同樣地，如令 $ \psi_3'=\psi_3+\alpha\psi_2'+\beta\psi_1' $ 並且由 $ (\psi_3',\psi_1')=(\psi_3',\psi_2')=0 $ 兩個正交條件可以解出 $ \alpha = -(\psi_2',\psi_3) / (\psi_2',\psi_2') $ 及 $ \beta = -(\psi_1',\psi_3)/ (\psi_1',\psi_1') $，並得到

\begin{align}
\psi_3'=\psi_3-\psi_2'\frac{(\psi_2',\psi_3)}{(\psi_2',\psi_2')}-\psi_1'\frac{(\psi_1',\psi_3)}{(\psi_1',\psi_1')}.
\end{align} 

從上式不難看出依序做下去可得到 $\psi_i'$ 

\begin{align}
\notag\psi_i'=\psi_i-\psi_{i-1}'\frac{(\psi_{i-1}',\psi_i)}{(\psi_{i-1}',\psi_{i-1}')}- \psi_{i-2}'\frac{(\psi_{i-2}',\psi_i)}{(\psi_{i-2}',\psi_{i-2}')}-\cdots-\psi_{1}'\frac{(\psi_{1}',\psi_i)}{(\psi_{1}',\psi_{1}')}.
\end{align}

$\{\psi_1',\psi_2',\ldots,\psi_n'\}$ 就是初步求得一組互相正交的新基底，即 $(\psi_i', \psi_j')=0$ 只要 $i\neq j$。第二個步驟是將這組正則化，如令


\begin{align}
\tilde{\psi_i}=\frac{\psi_i'}{\sqrt{(\psi_i',\psi_i')}},
\end{align}

\noindent 那麼 $\{\tilde{\psi_1},\tilde{\psi_2},\ldots,\tilde{\psi_n}\}$ 就是該 $n$ 維向量空間的一組正則基底，即 $\tilde{\psi_i}$ 與 $\tilde{\psi_j}$ 滿足 $(\tilde{\psi_i}, \tilde{\psi_j})=\delta_{ij}$。任何一個向量 $\psi$ 就可以由 $\psi_i$ 的線性組合來建構；即由 $n$ 個係數來組合，$\psi=\sum_{i=1}^n\alpha_i\psi_i$，係數 $\alpha_i=(\psi_i, \psi)$。今後當我們討論 $n$ 維向量空間的基底時，即使不特別說明，它們永遠代表的是一組正則基底。

\subsubsection*{範例六}

考慮 $ \mbox{C}(-\infty,\infty) $ 間的連續函數


\begin{align}
f_n(x)=x^n\exp\left(-\frac{x^2}{2}\right),~~~n=0,1,\ldots.
\end{align}

這組函數是 $ \mathcal{L}^2(-\infty,\infty)$ 空間的一組線性獨立向量，當然它們不一定互相正交。當 $n$ 與 $m$ 皆為偶數或皆為奇數時 $(f_n(x), f_m(x))\neq 0$，只有 $n$ 與 $m$ 一為奇數、一為偶數的時候 $f_n(x)$ 與 $f_m(x)$ 正交，即 $(f_n(x), f_m(x))= 0$。我們當然可以依 Gram-Schmidt 正交化的步驟來處理這組基底，先定義

\begin{align}
f_0'(x)&=f_0(x)=x^0\exp\left(-\frac{x^2}{2}\right),
\end{align}

然後建構 $f_1'(x)=f_1(x)+\alpha f_0'(x)$ 且令 $(f_1'(x), f_0'(x))=0$ 而求出 $\alpha=-(f_0'(x),f_1(x))/$ $(f_0'(x),f_0'(x)) = 0$，即 $f_1'(x)=x\exp\left(-x^2/2\right)$。再令 $f_2'(x)=f_2(x)+\alpha f_1'(x)+\beta f_0'(x)$，由於 $f_2'(x)$ 分別與 $f_1'(x)$ 和 $f_0'(x)$ 正交而可求得 $\alpha=0,~ \beta=-1/2$，換言之 $f_2'(x)=(x^2-1\2)\exp\left(-x^2/2\right)$。透過相同的步驟我們可以計算出下一個基底向量 $f_3'(x)=(x^3-3x/2 )\exp\left(-x^2/2\right)$。我們發現這一組新的基底連續函數 $f_n'(x)$ 可以表示成 $\exp\left(-x^2/2\right)$ 與 $x$ 的多項式的乘積，而這個特別的多項式 $H_n(x)$又稱為 Hermite 多項式，即

\begin{align}
\notag f_0'(x)&=2^0H_0(x)\exp\left(-\frac{x^2}{2}\right)=\exp\left(-\frac{x^2}{2}\right),\\
\notag f_1'(x)&=2^1H_1(x)\exp\left(-\frac{x^2}{2}\right)=2^1x\exp\left(-\frac{x^2}{2}\right),\\
\notag f_2'(x)&=2^2H_2(x)\exp\left(-\frac{x^2}{2}\right)=2^2\left(x^2-\frac{1}{2}\right)\exp\left(-\frac{x^2}{2}\right),\\
\notag \cdots\\
f_n'(x)&=2^nH_n(x)\exp\left(-\frac{x^2}{2}\right).
\end{align}

Hermite 多項式的性質將於下一章再仔細分析。如果我們再進一步將 $f_n'(x)$ 正則化，所得到的 $\tilde{f_n}(x)$ 則表為


\begin{align}
\tilde{f_n}(x)=\frac{f_n'(x)}{\sqrt{(f_n'(x),f_n'(x))}}=\frac{1}{\sqrt{2^nn!\sqrt{\pi}}}\exp\left(-\frac{x^2}{2}\right)H_n(x),
\end{align}


讀者不難認出 $f_n'(x)$ 其實就是一維簡諧振盪量子系統所解出來的本徵函數（eigenfunctions）。

\section{完備性與 Hilbert 空間}

向量空間裡，一個向量 $\psi $ 的範數（norm）以 $ \norm{\psi} $ 表之，它是一個非負值的函數；這裡所謂的函數，指的是向量 $psi$ 的函數。範數必須滿足下面三個條件：

\addtocounter{equation}{1}
\begin{align*}
&\norm{\psi}\geqslant 0,~~~~~ \mbox{當 } \psi \mbox{等於零向量時，即 } \psi=0，\mbox{則} \norm{\psi}=0 \tag{\arabic{chapter}.\arabic{equation}a}\\  
&\norm{a\psi}=|a|\cdot\norm{\psi},\tag{\arabic{chapter}.\arabic{equation}b}\\
&\norm{\psi +\phi}\leqslant \norm{\psi}+\norm{\phi}.\tag{\arabic{chapter}.\arabic{equation}c} 
\end{align*}

向量的內積和向量的範數（norm）各有不同的意義及其必須滿足的條件。但是物理學者常把這兩者混為一談，原因在於 Schwarz 不等式可以透過內積操作導出，因此如果把 $\norm{\psi}^2$ 認定為 $(\psi, \psi)$，那麼向量的內積也可以滿足前述範數定義中三個條件，以致兩者混合使用。嚴格來說，範數的條件比內積來的寬鬆，我們會以範例七來說明。

\subsubsection*{範例七}

令 $ f(x)$ 為  $\mbox{C}(a,b) $ 的函數，即 $ f(x) $ 為 $ a \leqslant x\leqslant b $ 的連續函數。我們可以定義 $\norm{f(x)}$ 為 $\sqrt{((f(x),f(x))}$ ，因為

\addtocounter{equation}{1}
\begin{align*}
(f(x),f(x))&=0,\mbox{代表 } f(x) {在 }a \leqslant x\leqslant b \mbox{間恆為零}. \tag{\arabic{chapter}.\arabic{equation}a}\\  
\norm{af(x)}^2&=(af(x), af(x))=|a|^2(f(x), f(x))\tag{\arabic{chapter}.\arabic{equation}b}\\
(\norm{\psi}+\norm{\varphi})^2&=\norm{\psi}^2+2\norm{\psi}\cdot\norm{\varphi}+\norm{\varphi}^2\\
& \geqslant \norm{\psi}^2+2|\psi\cdot\varphi |+\norm{\varphi}^2\geqslant \norm{\psi+\varphi}^2.\tag{\arabic{chapter}.\arabic{equation}c} 
\end{align*}

\subsection{Cauchy 序列 and convergent sequence}
Consider an infinite dimensional vector space and denote the basis by $ \{\phi_1,\phi_2,\phi_3,\ldots \} $. We construct the partial sum $ \psi_N=\sum_ia_i\phi_i $, where $ i $ runs from $ 1 $ to $ N $, and obtain $ \ldots,\psi_j,\psi_{j+1},\ldots ,\psi_m,\psi_{m+1},\ldots, \psi_n,\ldots $ for increasing values in $ N $ that forms an infinite sequence. The sequence is called a \textbf{Cauchy sequence}\index{Cauchy sequence} if 
$$
\lim_{n\rightarrow \infty} \psi_n=\lim_{m\rightarrow \infty} \psi_m,~~~~~~~~ 
$$
or more precisely to put in terms of norm, i.e, $ \displaystyle\lim_{n,m\rightarrow \infty}\norm{\psi_n-\psi_m}=0 $. 

It is said that a vector $ \psi_m $ converges to $ \psi $ if 

$$
\lim_{m\rightarrow\infty}\psi_m =\psi,  ~~~\mbox{or}~~~ \lim_{m\rightarrow \infty}\norm{\psi_m-\psi}=0, 
$$ 

\noindent then $ \{\ldots,\psi_{m-1},\psi_{m},\ldots\} $ is called a \textbf{convergent sequence}\index{convergent sequence}.

It is easily concluded that every convergent sequence is a Cauchy sequence. Yet it is not necessary true conversely. Namely a Cauchy sequence is not always a convergent sequence.

\subsection{Complete vector space}
\index{completeness!complete vector space}\index{space!complete vector space}
A vector space, in which every Cauchy sequence of a vector $ \psi_m $ converges to a limiting vector $ \psi $, is called a complete vector space.

\subsection{Hilbert space}
\index{Hilbert space}\index{space!Hilbert space}
A Hilbert space is a complete vector space with norm defined as the inner product. A Hilbert space, finite dimensional or infinite dimensional, is separable if its basis is countable.
\chapter{量子力學的各種表示理論}

\section{有限維與無限維的向量空間}

在第一章討論向量空間時，我們一直採以抽象的符號來說明向量與向量空間的種種性質與相關定理。從這一章開始，我們將以學習過的數學，或者新介紹的數學工具，重新驗證上一章以抽象符號描述的理論。

首先我們以 $\mathcal{C}^n$ 空間為例，過去我們學習過的矩陣理論，其實就是有限維向量空間中最具體的表示（representation）。在 $\mathcal{C}^n$ 空間中的一個向量，通常用 $n$ 個複數由上而下排成一行來表示，譬如抽象的向量 $ \bar{x} $ 即可表示為

\begin{align}
\bar{x}=\begin{pmatrix}x_1\\x_2\\\vdots\\x_i\\x_{i+1}\\\vdots\\x_n\end{pmatrix},
\end{align}

\noindent 這個向量矩陣我們稱為行矩陣（column matrix），其中 $x_i$ 是該向量的第 $i$ 個分量，當然是複數。因為矩陣的運算操作法則與向量的兩個基本操作法則一致，並且也滿足向量的八個條件，這就是為何數學家以矩陣來進行有限維向量空間的各種運算。

用行矩陣來表示 $\mathcal{C}^n$ 空間中的一個向量並不是唯一的選擇或方法。我們也曾在第一章時，將一個向量 $x$ 用橫寫得 $n$ 個 $x_i$，即 $(x_1, x_2, \ldots, x_i, \ldots, x_n)$ 來表示。也就是說 $n$ 維向量也可以用列矩陣（row matrix） 來描述。用列矩陣來表示向量的線性空間稱為$\mathcal{C}^n$ 空間的對偶空間（dual space），使用符號 $ \tilde{\mathcal{C}}^n $ 來表示。為了有別於前述的行矩陣 $\bar{x}$，$ \tilde{\mathcal{C}}^n $ 中的向量符號為：

\begin{align}
\tilde{x}=\begin{pmatrix}x_1,&x_2,&\ldots, & x_i, & x_{i+1}, & \ldots, & x_n \end{pmatrix},
\end{align}

矩陣的表示法有很多方便之處。$\mathcal{C}^n$ 空間的第 $i$ 個基底向量就很自然地表為 $\bar{e_i}$，它也是具有 $n$ 個分量的行矩陣，除了第 $i$ 個分量為 1，其他分量皆為零，即

\begin{align}\label{columnmatrix}
\bar{e_i}=\left(\begin{array}{c}0\\0\\\vdots\\1\\0\\\vdots\\0\end{array}\right),
\end{align}

任何一個向量 $\bar{x}$ 即可表以該組基底的線性組合，即

\begin{align}
\bar{x}=\sum_{i=1}^nx_i\bar{e_i}.
\end{align}

在 $ \tilde{\mathcal{C}}^n $ 空間中，第 $i$ 個基底向量則寫為

\begin{align}
\tilde{e_i}=\begin{pmatrix}0,&0,&\ldots, & 1, & \ldots, & 0 \end{pmatrix}.
\end{align}

所以任何 $ ^*\mathcal{C}^n $ 空間中的向量 $\tilde{x}$ 也可以表示成 $\tilde{e_i}$ 的線性組合，即 

\begin{align}
\tilde{x}=\sum_{i=1}^nx_i\tilde{e_i}.
\end{align}

以矩陣討論量子力學時， $\tilde{x}$ 空間的觀念用處不多，反而另一個對偶空間 $ ^*\mathcal{C}^n $ 的用處更大。 $ ^*\mathcal{C}^n $ 空間中的向量也是一個列矩陣，在 $ \mathcal{C}^n $ 中的一個向量 $\bar{x}$ 對應於 $ ^*\mathcal{C}^n $ 中的向量寫為 $\tilde{x}^*$，這個向量是行矩陣 $\bar{x}$ 的伴隨共軛形（adjoint conjugate），亦即將 $\bar{x}$ 的各個分量 $x_i$ 的共軛（complex conjugate）依序擺入列矩陣中：

\begin{align}
\tilde{x}^*=\bar{x}^\dag=\begin{pmatrix}x_1^*,&x_2^*,&\ldots, & x_i^*, & x_{i+1},^* & \ldots, & x_n^* \end{pmatrix}
\end{align}

$ ^*\mathcal{C}^n $ 空間的好處可以從向量空間的內積計算看出來。當我們求取 $\bar{x}$ 與 $\bar{y}$ 這組有序對的內積時我們就可以遵循矩陣的基本操作法則，計算 $\bar{x}$ 向量的伴隨共軛 $\bar{x}^\dag$ 與 $\bar{y}$ 向量的積。也就是內積 $(x,y)$ 可寫為列矩陣 $\tilde{x}^*$ 與行矩陣 $\bar{y}$ 的積。

\begin{align}
(x,y)=\tilde{x}^*\bar{y}=\bar{x}^\dag\bar{y}=\begin{pmatrix}x_1^*&x_2^*&\ldots & x_i^*&\ldots &x_n^*\end{pmatrix}\begin{pmatrix}y_1\\y_2\\\vdots \\y_i\\\vdots\\y_n\end{pmatrix}=\sum_{i=1}^nx_i^*y_i.
\end{align}

值得一提的是 $ \tilde{\mathcal{C}}^n $ 空間及 $ ^*\mathcal{C}^n $ 空間的基底構成相同，第 $i$ 個基底向量都可寫為

\begin{align}\label{rowmatrix}
\tilde{e_i}^*=\tilde{e_i}=\begin{pmatrix}0,&0,&\ldots, & 1, & \ldots, & 0 \end{pmatrix}.
\end{align}

由 eq.~\ref{rowmatrix} 與 eq.~\ref{columnmatrix} 可以得到


\begin{align}
\tilde{e_i}\bar{e_j}=\delta_{ij}.
\end{align

\noindent 上式說明了 $(\bar{e_1),\bar{e_2),\ldots, \bar{e_i), \ldots, \bar{e_n))$ 這 $n$ 個 $ \mathcal{C}^n $ 空間的基底已經正則化（orthonormalized）。

$\mathcal{C}^n $ 空間中線性算子（linear operator）在矩陣的表示法就是大家所熟悉的 $n\times n$ 矩陣。如果我們用相同的符號 $ \mathbf{A} $ 來表示 $n\times n$ 矩陣，則矩陣元素（matrix elements）$A_{ij}$ 可寫為：

\begin{align}\label{matrix-present1}
A_{ij}=\tilde{e_i}\mathbf{A}\bar{e_j}.
\end{align}

因為 $\bar{e_j}$ 的伴隨共軛為 $\bar{e_j}$，$\tilde{e_i}$ 的伴隨共軛為 $\bar{e_i}$，所以 $ \mathbf{A}^\dag $ 的矩陣元素 $(A_{ij})^\dag$ 為

\begin{align}
(A_{ij})^\dag=\tilde{e_i}\mathbf{A}^\dag\bar{e_j}=\bar{e_i}^\dag A^\dag \tilde{e_j}^\dag=(\tilde{e_j}\mathbf{A}\bar{e_i})^\dag.
\end{align}

但 $\tilde{e_j}\mathbf{A}\bar{e_i}$ 唯一純數，所以它的伴隨形（adjoint）就是取複共軛（complex conjugate），因此

\begin{align}
(A^\dag)_{ij}=(\tilde{e_j}\mathbf{A}\bar{e_i})^\dag=(\tilde{e_j}\mathbf{A}\bar{e_i})^*
= A_{ji}^*.
\end{align}

\noindent 驗證了 $\mathbf{A}$ 的伴隨共軛形式。

如果 $\mathbf{A}$ 代表量子系統中的可觀測量（observable），那麼 $\mathbf{A}=\mathbf{A}^\dag$，當該系統在量子態 $\bar{x}$ 進行 $\mathbf{A}$ 的測量時，它的期望值可表以 $\bar{x}^\dag\mathbf{A}\bar{x}$ 三個矩陣的乘積：

\begin{align}
\expecto{\mathbf{A}}=\bar{x}^\dag\mathbf{A}\bar{x}=\bar{x}^\dag\mathbf{A}^\dag\bar{x} =(\bar{x}^\dag\mathbf{A}\bar{x})^\dag=(\bar{x}^\dag\mathbf{A}\bar{x})^*=\expecto{\mathbf{A}}^*.
\end{align} 

因此我們再次用矩陣來驗證量子力學中的可觀測量之期望值是實數。

有限維向量的矩陣表示法對於建構投影算子（projection operator）極為方便。如果我們把基底向量（行矩陣） $\bar{e_i}$ 寫在左邊，並在右邊乘上共軛空間的基底向量（列矩陣） $\tilde{e_j}$，則該矩陣乘積就是一個 $i$ 維子空間（$i$ subspace）的投影算子 $\mathbf{P}_i$，即


\begin{align}
\mathbf{P}_i=\bar{e_i}\tilde{e_i}=\left(\begin{array}{c}0\\0\\\vdots\\1\\0\\\vdots\\0\end{array}\right)\begin{pmatrix}0&0&\ldots & 1 & 0 & \ldots & 0 \end{pmatrix}=
\begin{pmatrix}
0&0&\ldots&0&0&\ldots&0\\
0&\ddots&&&&&\vdots\\
\vdots&&\ddots&&&&\vdots\\
0&&&1&&&0\\
0&&&&0&&0\\
\vdots&&&&&\ddots&\vdots\\
0&0&\ldots&0&0&\ldots&0\\
\end{pmatrix}.
\end{align}

讀者可自行驗證

\begin{align*}
\mathbf{P}_i\mathbf{P}_j=\left\{\begin{array}{l}0, \mbox{當}~ i\neq j, \\ \mathbf{P}_i, \mbox{當}~ i\neq j, ~\mbox{亦即}~ \mathbf{P}_i^2=\mathbf{P}_i.\end{array}
\end{align*}

一個可以投影到子空間 $\mathcal{M}$ 的投影算子 $\mathbf{P}_\mathcal{M}$ 則為 $\mathbf{P}_i$ 的部份和，亦即將 $\mathcal{M}$ 中的所有 $e_i$ 的投影算子相加：


\begin{align}
\mathbf{P}_\mathcal{M}=\sum_{ e_i\in\mathcal{M}}\mathbf{P}_i.
\end{align}

如果 $\mathcal{M}$ 取為 $n$ 維的全部空間，則


\begin{align}
\sum_{i=1}^n\mathbf{P}_i=\sum_{i=1}^n\bar{e_i}\tilde{e_j}=\mathbf{I}.
\end{align}

\noindent 上式就是封閉性質（closure relation）的矩陣表達。

對應於投射在子空間中向量 $\bar{y}$ 的投影算子 $\mathbf{P}_{\bar{y}}$ ，我們也可以利用行矩陣 $\bar{y}$ 與其伴隨共軛的列矩陣 $\tilde{y}$ 有序操作來建構出投影算子：

\begin{align}
\mathbf{P}_y=\bar{y}\bar{y}^\dag=\bar{y}\tilde{y}^*.
\end{align}

矩陣表示法適用於有限維的向量空間，它無法應用在無限維空間的討論與分析。但是量子力學涉及的 Hilbert 空間往往都是無限維度，並且是不可分離的（nonseparable）。幸好在 1930 年代，Dirac  設計了一套符號適合用來討論及分析有限維或無限維、可分離或者不可分離的 Hilbert 空間，直至今日這套簡潔的符號系統仍廣為物理學家使用。

在 Hilbert 空間 $\mathcal{H}$ 中的一個向量， Dirac 用符號 $\ket{\alpha}$ 來表示，並稱之為 ket 向量。對應於 Hilbert 空間 $\mathcal{H}$ 的對偶空間 $^*\mathcal{H}$ 的伴隨共軛向量則使用 符號 $\bra{\alpha}$  來註記，並稱為 bra 向量。當我們要計算一對有序向量 $\beta$ 與 $\alpha$ 的內積時，可以用 Dirac 的符號系統（bra-ket）將 $(\beta, \alpha)$ 寫為 $\bra{\alpha}\cdot\ket{\beta}$， 或者簡寫為 $\braket{\beta}{\alpha}$，也就是從 $^*\mathcal{H}$ 取一個 bra 向量 $\bra{\alpha}$ 和 $\mathcal{H}$ 空間中的 ket 向量 $\alpha$ 依照有序對（order pair）的規則乘在一起：

\begin{align}
\bra{\alpha}\ket{\beta}\equiv\braket{\alpha}{\beta}.
\end{align}

當 $\beta$ 取為 $\alpha$ 時，相當於得到 $\mathcal{H}$ 的範數平方，

\begin{align}
\norm{\alpha}^2=\braket{\alpha}{\alpha}.
\end{align}

如果 $\mathbf{A}$ 代表 $\mathcal{H}$ 空間中的一個可觀測量，當 $\mathbf{A}$ 作用於 $\ket{\alpha}$ 而得到 $ \ket{\alpha'} $，即 $\mathbf{A}\ket{\alpha}=\ket{\alpha'}$ ，我們計算 $\ket{\beta}$ 與 $\ket{\alpha'}$ 內積時可寫為

\begin{align}
(\beta, \alpha')=\braket{\beta}{\alpha'} = \bra{\beta}\left(\mathbf{A}\ket{\alpha}\right) \equiv \bra{\beta}\mathbf{A}\ket{\beta}.
\end{align}

上式中間項的括號往往可以省略，並且意味著算子 $\mathbf{A}$ 是作用於右邊的 ket 向量。如果計算 $\ket{beta'}=\mathbf{A}\ket{\beta}$ 和 $\ket{\alpha}$ 的內積時，則因為 $\bra{\beta'}=(\ket{\beta'})^\dag=\bra{\beta}\mathbf{A}^\dag$，所以

\begin{align}
(\beta',\alpha)=(\mathbf{A}\beta, \alpha)=(\bra{\alpha}\mathbf{A}^\dag)\ket{\beta}=\bra{\alpha}\mathbf{A}^\dag\ket{\beta}.
\end{align}

上式第三項的括號也可省略。這表示該內積

\begin{align}
\bra{\beta}\mathbf{A}^\dag\ket{\alpha}=\bra{\beta}(\mathbf{A}^\dag\ket{\alpha})=(\beta, \mathbf{A}^\dag\alpha),
\end{align}

可寫為

\begin{align}
(\beta', \alpha)=(\mathbf{A}\beta, \alpha)=(\beta, \mathbf{A}^\dag\alpha).
\end{align}

該是即為上一章介紹伴隨算子（adjoint operator）$\mathbf{A}^\dag$ 的定義。

如果 $\mathbf{A}$ 代表的是一個量子系統的可觀測量時，則 $\mathbf{A}=\mathbf{A}^\dag$，當 $\beta$ 取為 $\alpha$ 時，$(\bra{\alpha}\mathbf{A}^\dag\ket{\alpha})=(\bra{\alpha}\mathbf{A}\ket{\alpha})^\dag= \bra{\alpha}\mathbf{A}^\dag\ket{\alpha}^*$，所以量子可觀測量的期望值為實數亦可由 Dirac 符號系統再次驗證。

第一章我們所討論的量子力學定律，如果以 Dirac 符號來描述，則顯得更為簡潔。我們以量子系統的算子 $\mathbf{A}=\mathbf{A}^\dag$ 來建構一個 $\mathcal{H}$ 空間，那麼該空間的基底就可從 $\mathbf{A}\ket{a}=a\ket{a}$ 的固有向量（eigenvecotr，字根 eigen- 為德語『具備某種特徵』、『特定於某物』之意）方程式解出所有的 $\ket{a}$ 而獲得。當固有值（eigenvalue） $a$ 是離散值，我們可用 $a_1, a_2, a_3, \ldots, a_i, \ldots $ 來表示，當足碼 $i$ 由 1 至 $n$ 就能描述所有可能的 $a_n$ 時，那麼該 Hilbert 空間就是 $n$ 維度空間。如果 $a_i$ 可以取無限多值，那麼這個 $\mathcal{H}$ 空間為無限多維度可分離的 Hilbert 空間。任何該系統的量子狀態 $\ket{\alpha}$ 就可表以 $\ket{\alpha_i}$ 的線性組合，即

\begin{align}
{\alpha} = \sum_j\alpha_j\ket{a_j}.
\end{align}

當 $\ket{\alpha_i}$ 已正則化時，

\begin{align}
\braket{a_i}{a_j}=\delta_{ij},
\end{align}

那麼 $\alpha_i$ 可從內積 $\braket{a_i}{\alpha}$ 求出

\begin{align}
\alpha_i=\braket{a_i}{\alpha}=\bra{a_i}\left(\sum_j\alpha_j\ket{a_j}\right)=\sum_j\alpha_j\delta_{ji}.
\end{align}

我們可以獲得 Dirac 符號所表示的封閉性質

\begin{align}
\sum\proj{a_i}{a_i}=\mathbf{I}.
\end{align}

其中 $\mathbf{P}_i=\proj{a_i}{a_i}$ 就是投射於第 $i$ 個子空間的投影算子。如果是投射於子空間 $\mathcal{M}$ ，則

\begin{align}
\mathbf{P}_\mathcal{M}=\sum_{i}\mathbf{P}_i=\sum_i\proj{a_i}{a_i},~~~\ket{a_i}\in\mathcal{M}.
\end{align}

當我們解固有向量方程式 $\mathbf{A}\ket{a_i}=a_i\ket{a_i}$，如果對應於 $a_i$ 具有 $m$-folde 的簡併（degeneracy）時，則該固有向量必須多一個足碼來加以區別簡併態，我們通常將它表以 $\ket{a_i,l}$ 其中 $a_i$ 為固有值，$l=1,2,\ldots, m$ 代表該固有向量的 $m$ 個簡併態。該系統的任何一個量子態 $\ket{alpha}$ 則表為

 
\begin{align}
\ket{\alpha}=\sum_i\sum_{l=1}^m\alpha_{i,l}\ket{\alpha_{i,l}},
\end{align}

其基底 $\ket{\alpha_{i,l}}$ 的正則化寫為

\begin{align}
\braket{a_{i,l}}{a_{j,k}}=\delta_{ij}\delta_{jk}.
\end{align}

量子系統中，對應於某些可觀測量所求出的固有值往往是連續而非離散的。我們也可以用它的固有值來標示其對應的固有向量，比如

\begin{align}
\mathbf{A}\ket{a}=a\ket{a}.
\end{align}

$a$ 是一個連續值的實數，因此不可分離的 $\mathcal{H}$ 是無限維的。固有向量 $\ket{a}$ 的正則化當然有別於上述離散的情況，我們必須引入一個新的符號： $\delta$ 函數來進行描述，將 $\ket{a}$ 及 $\ket{a'}$ 的內積寫為

\begin{align}
\braket{a}{a'}=\delta(a-a').
\end{align}

因為 $a\neq a'$ 時，$\ket{a}$ 與 $\ket{a'}$  正交，所以 $\delta(a-a')=0$>。但是 $a'$ 可無限接近 $a$ ，所以相似於可分離的固有向量的情況，即 $\braket{a_i}{a_j}=\delta_{ij}$ 且 $\sum_i\braket{a_i}{a_j}=\sun_i\delta{ij}=1$，我們可以定義

\begin{align}
\int\braket{a}{a'} d\a=\int\deta{a-a')=1.
\end{align}

也就是說，$\delta(x-a)$ 函數具有以下性質：

\begin{align*}
\delta(x-a)\left\{\begin{array}{l}=0,\mbox{~當}~x\neq a ~\mbox{時},\\\mbox{未定，當}~x= a ~\mbox{時}.\end{array}
\end{align*}

但 $\delta(x-a)$ 的總面積為 1，即 $a$ 在積分範圍 $\Omega$ 之內時

\begin{align*}
\int{\Omega}\delta(x-a)dx=1.
\end{align*}

結束本節之前，我們還要在介紹一個經常使用，不可分離的 $\mathbf{H}$ 的封閉性質。它也是從 $\sum_i\proj{i}{i}=\mathbf{I}$ 的離散和延伸到 $\ptoj{a}{a}$ 對 $a$ 的積分，即

\begin{align*}
\int\ket{a}da\bra{a}=\int \mathbf{P}_ada=\mathbf{I},
\end{align*}

上式中的 $\mathbf{P}_a$ 式投影到基底 $\ket{a}$ 的投影算子 $\mathbf{P}_a=\ket{a}\bra{a}$。

\section{量子力學的 q 表示與 p 表示}

量子系統的實際演算，最常用的是 q 表示（q-representation）。因為大多數的讀者最熟悉的數學也是微積分，所以操作微分算子以及波函數要比操作抽象的線性算子或者量子態來得自然。這一節我們要討論如何將抽象的 $\mathcal{H}$ 空間的量子系統以讀者熟悉的 q 表示來呈現計算的細節。

為了簡要起見，我們只以一維的空間作為示範。首先我們選取對應於位置的量子可觀測量 $ \mathbf{X} $ ，並由 $ \mathbf{X} $ 的固有向量 $\ket{x}$ 來展（span）成一個無限維的 $\mathcal{H}$ 空間，即由

\begin{align}
\mathbf{X}\ket{x}=x\ket{x}
\end{align}

\noindent 所獲得的不可分離固有向量 $\ket{x}$ 作為基底來建構一個 $\mathcal{H}$ 空間。如果我們將該量子系統進行位移，例如由 $x$ 的位置移到 $x+\xi$ 的位置，也就是說，我們試著建構一個平移算子（translational operator），以 $\mathbf{U}(\mathbf{P};\xi)=\exp( -\frac{i}{\hbar}\xi\mathbf{P})$ 的符號來表示，那麼將算子作用到固有向量上會得到

\begin{align}
\mathbf{U}(\mathbf{P};\xi)\ket{x}=e^{-\frac{i}{\hbar}\xi\mathbf{P}}=\ket{x＋\xi}.
\end{align}

如何求得該平移算子 $\mathbf{U}(\mathbf{P};\xi)$ 為 $\exp( -\frac{i}{\hbar}\xi\mathbf{P})$ 會留待下一章對稱轉換一節中加以討論，我們在此先討論該算子的特性：

\begin{align*}
\xi&=\mbox{實參數, 即空間中平移的距離}\\
\mathbf{P}&=\mbox{一維的動量算子},
\end{align*}

並且 

\begin{align*}
\mathbf{U}^\dag(\mathbf{P};\xi) = e^{\frac{i}{\hbar}\xi\mathbf{P}}= \mathbf{U}^{-1}(\mathbf{P};\xi).
\end{align*}  

所以 $\mathbf{U}(\mathbf{P};\xi)$ 是一個么正算子（unitary operator）。我們先證明第一個引理：

\proposition{
\subsubsection*{引理一}
$ \mathbf{U}(\mathbf{P};\xi)\ket{x} $ 為位置算子 \mathbf{X} 的固有向量，其固有值為 $x+\xi$; 即 $\mathbf{U}(\mathbf{P};\xi)\ket{x}=\ket{\alpha}$，則 $\mathbf{X}\ket{\alpha}=(x+\xi)\ket{\alpha}$，所以 $\ket{\alpha}$ 可寫為 

\begin{align*}
\ket{\alpha}=\mathbf{U}(\mathbf{P};\xi)\ket{x}=\ket{x+\xi}.
\end{align*}}

\subsection*{證明}

從基本對易關係 $ [\mathbf{X},\mathbf{P}]=i\hbar\mathbf{I} $ 我們不難計算 $[\mathbf{X},\mathbf{P}]^n$ 如下：

\begin{align}
[\mathbf{X},\mathbf{P}^n]=\mathbf{P}^{n-1}[\mathbf{X},\mathbf{P}]+[\mathbf{X},\mathbf{P}^{n-1}]\mathbf{P}=in\hbar\mathbf{P}^{n-1},
\end{align}

故可進一步得到

\begin{align}
\notag [\mathbf{X},\mathbf{U}(\mathbf{P};\xi)]&=i\hbar\left[\mathbf{X},\sum_m\frac{1}{m!}\left(-\frac{i}{\hbar}\xi\right)^m\mathbf{P}^m\right]\\
&=\xi\sum_m\frac{1}{(m-1)!}\left(-\frac{i}{\hbar}\xi\right)^{m-1}\mathbf{P}^{m-1}=\xi\mathbf{U},
\end{align}

或寫為

\begin{align}
\mathbf{XU}(\mathbf{P};\xi)=\mathbf{U}(\mathbf{P};\xi)(\mathbf{X}+\xi\mathbf{I}).
\end{align}

當 $\mathbf{U}(\mathbf{P};\xi)\ket{x}=\ket{\alpha}$ 時，則

\begin{align}
\mathbf{X}\ket{\alpha}=\mathbf{XU}\ket{x}=\mathbf{U}(\mathbf{X}+\xi\mathbf{I})\ket{x}=\mathbf{U}(x+\xi)\ket{x}=(x+\xi)\ket{\alpha},
\end{align}

上式表示 $\ket{\alpha}$ 也是 $\mathbf{X}$ 的固有向量，其固有值為 $x+\xi$，亦即 $\ket{\alpha}=c\ket{x+\xi}$。通常 $c$ 為相位因子。為了簡便，通常把 $c$ 取為 1，所以引理一證實了

\begin{align}
\mathbf{U}(\mathbf{P};\xi)\ket{x}=\ket{x+\xi}.
\end{align}

\end{document}